% !TeX program = xelatex
\documentclass[11pt,a4paper]{article}
\usepackage[margin=22mm,headheight=14pt,headsep=8mm,footskip=10mm]{geometry}
\usepackage{fontspec}
\setmainfont{Times New Roman}
\setsansfont{Arial}
\setmonofont{Consolas}[Scale=0.88]
\usepackage{amsmath,amssymb,graphicx,booktabs,longtable,tabularx,array}
\usepackage{enumitem,fancyhdr,titlesec,caption,needspace,fvextra,xurl,float}
\usepackage[hidelinks,unicode,pdftitle={Hệ thống trợ lý tư vấn sách ứng dụng nhiều tác tử},pdfauthor={Hoàng Xuân Thanh, Trần Gia Lạc}]{hyperref}
\usepackage{tikz}
\hypersetup{bookmarksnumbered=true}
\renewcommand{\figurename}{Hình}
\renewcommand{\tablename}{Bảng}
\renewcommand{\contentsname}{Mục lục}
\renewcommand{\refname}{Tài liệu tham chiếu}
\captionsetup{font=small,labelfont=bf,justification=raggedright,singlelinecheck=false}
\setlength{\parindent}{0pt}
\setlength{\parskip}{5.5pt}
\linespread{1.09}
\setlist{nosep,leftmargin=5mm,itemsep=3pt,topsep=4pt}
\setlength{\emergencystretch}{2.5em}
\renewcommand{\arraystretch}{1.16}
\newcolumntype{Y}{>{\raggedright\arraybackslash}X}
\newcolumntype{P}[1]{>{\raggedright\arraybackslash}p{#1}}
\titleformat{\section}{\normalfont\large\bfseries}{\thesection.}{0.6em}{}
\titleformat{\subsection}{\normalfont\normalsize\bfseries}{\thesubsection.}{0.6em}{}
\titlespacing*{\section}{0pt}{13pt}{6pt}
\titlespacing*{\subsection}{0pt}{9pt}{4pt}
\newcommand{\src}[1]{\hyperlink{src-#1}{[#1]}}
\ExplSyntaxOn
\NewDocumentCommand{\code}{m}{
  \tl_if_in:nnTF {#1} {~}
    {\texttt{\detokenize{#1}}}
    {\nolinkurl{#1}}
}
\ExplSyntaxOff
\newcommand{\file}[1]{\nolinkurl{#1}}
\newcommand{\fig}[3]{\begin{figure}[H]\centering\includegraphics[width=\linewidth]{images/#1.png}\caption{#2}\label{#3}\end{figure}}
\pagestyle{fancy}
\fancyhf{}
\fancyhead[L]{\small Hệ thống trợ lý tư vấn sách}
\fancyhead[R]{\small Dự án Công nghệ thông tin}
\fancyfoot[C]{\small\thepage}
\renewcommand{\headrulewidth}{0.3pt}
\widowpenalty=10000
\clubpenalty=10000
\begin{document}
% ---------- Trang bìa ----------
\begin{titlepage}
    \thispagestyle{empty}
    \centering

    \begin{tikzpicture}[remember picture,overlay]
        \draw[line width=0.6pt]
        ([xshift=1cm,yshift=-1cm]current page.north west)
        rectangle
        ([xshift=-1cm,yshift=1cm]current page.south east);
    \end{tikzpicture}

    \vspace*{0.2cm}

    {\large\bfseries TỔNG LIÊN ĐOÀN LAO ĐỘNG VIỆT NAM\par}
    \vspace{0.2cm}
    {\large\bfseries TRƯỜNG ĐẠI HỌC TÔN ĐỨC THẮNG\par}
    \vspace{0.2cm}
    {\large\bfseries KHOA CÔNG NGHỆ THÔNG TIN\par}

    \vspace{0.8cm}
    \includegraphics[width=4cm]{images/tdtu.jpg}\par

    \vspace{1cm}

    {\normalsize\bfseries HOÀNG XUÂN THANH -- 523H0178\par}
    \vspace{0.2cm}
    {\normalsize\bfseries TRẦN GIA LẠC -- 523H0154\par}

    \vspace{1.2cm}

    {\Large\bfseries
    HỆ THỐNG TRỢ LÝ TƯ VẤN SÁCH\\
    ỨNG DỤNG NHIỀU TÁC TỬ\par}

    \vspace{1.1cm}

    {\Large\bfseries BÁO CÁO MÔN HỌC\par}
    \vspace{0.25cm}
    {\Large\bfseries DỰ ÁN CÔNG NGHỆ THÔNG TIN\par}

    \vspace{1cm}

    {\normalsize\bfseries Giảng viên\par}
    \vspace{0.2cm}
    {\normalsize\bfseries PGS. TS. Lê Anh Cường\par}

    \vfill

    {\large\bfseries THÀNH PHỐ HỒ CHÍ MINH, 2026\par}
    \vspace*{0.2cm}
\end{titlepage}

% ---------- Lời cảm ơn ----------
\pagenumbering{roman}
\section*{LỜI CẢM ƠN}
\addcontentsline{toc}{section}{Lời cảm ơn}

Nhóm chúng em xin chân thành cảm ơn PGS. TS. Lê Anh Cường đã hướng dẫn, định hướng và góp ý trong quá trình thực hiện học phần Dự án Công nghệ thông tin. Những nhận xét của thầy giúp nhóm làm rõ bài toán, hoàn thiện cách trình bày hệ thống và nhìn nhận kết quả thực nghiệm thận trọng hơn.

Chúng em cũng trân trọng ghi nhận các bộ dữ liệu, tài liệu kỹ thuật và phần mềm mã nguồn mở được sử dụng trong quá trình xây dựng và đánh giá hệ thống. Các nguồn này được dẫn trong báo cáo để người đọc có thể tra cứu và đối chiếu.

Mặc dù đã cố gắng hoàn thiện đề tài, báo cáo vẫn có thể còn thiếu sót. Nhóm mong nhận được góp ý của thầy để tiếp tục cải thiện hệ thống và cách đánh giá trong thời gian tới.

\begin{flushright}
\textit{Thành phố Hồ Chí Minh, 2026}\\
\textbf{Nhóm thực hiện}\\[1cm]
Hoàng Xuân Thanh\\
Trần Gia Lạc
\end{flushright}

\clearpage

% ---------- Lời cam đoan ----------
\section*{LỜI CAM ĐOAN}
\addcontentsline{toc}{section}{Lời cam đoan}

Chúng em cam đoan báo cáo này được thực hiện trong khuôn khổ học phần Dự án Công nghệ thông tin. Nội dung mô tả thiết kế, dữ liệu, quy trình và kết quả được trình bày dựa trên tài liệu và hiện vật của dự án; các nguồn tham khảo được ghi nhận trong báo cáo.

Nhóm chịu trách nhiệm về tính trung thực của phần trình bày và việc ghi nguồn. Các kết quả đánh giá được diễn giải trong phạm vi bộ câu hỏi, dữ liệu và phương pháp chấm đã nêu; không được xem là bằng chứng cho mọi tình huống sử dụng thực tế.

\begin{flushright}
\textit{Thành phố Hồ Chí Minh, 2026}\\
\textbf{Nhóm thực hiện}\\[1cm]
Hoàng Xuân Thanh\\
Trần Gia Lạc
\end{flushright}

\clearpage

% ---------- Tóm tắt ----------
\section*{TÓM TẮT}
\addcontentsline{toc}{section}{Tóm tắt}

Đề tài xây dựng hệ thống trợ lý tư vấn sách tiếng Việt theo hướng phối hợp nhiều tác tử chuyên trách. Bộ phận điều phối chia yêu cầu thành các nhiệm vụ tra cứu thông tin sách, tổng hợp đánh giá người mua, nhận diện dấu hiệu phàn nàn và đối chiếu nguồn kiến thức về tác phẩm. Bước tổng hợp gắn nhận định với nguồn tương ứng và nêu rõ khi dữ liệu chưa đủ. Hệ thống sử dụng dữ liệu sách, đánh giá đã lưu và một tập nguồn được tuyển chọn; không cung cấp giá hay tình trạng tồn kho theo thời gian thực.

Báo cáo so sánh cách xử lý một tác tử và nhiều tác tử trên cùng 24 câu hỏi, mỗi câu được chạy ba lần, với cùng mô hình ngôn ngữ và thông tin liên quan đã cấp sẵn. Theo bộ chấm tự động dựa trên 43 dữ kiện bắt buộc, cách nhiều tác tử đạt 129/129 lượt dữ kiện (100\%), còn cách một tác tử đạt 126/129 (97,67\%). Thời gian trả lời trung bình lần lượt là 6,75 giây và 1,43 giây; số token trung bình là 3.706 và 1.332.

Phần chênh lệch điểm tập trung ở một câu hỏi về tên ấn bản. Vì vậy, kết quả hiện tại cho thấy một đánh đổi giữa mức đáp ứng dữ kiện và tài nguyên sử dụng, chưa đủ để kết luận nhiều tác tử ưu việt nói chung. Kết quả cũng phản ánh phép thử có thông tin liên quan cấp sẵn và bộ chấm tự động, không thay thế đánh giá của con người hay chứng minh chất lượng trong mọi tình huống thực tế.

Sau thực nghiệm, nhóm sửa sáu điểm khiến hệ thống có thể làm rơi dữ kiện và loại khỏi kết luận các số liệu được viết sẵn trong mã chương trình. Các sửa đổi được xác nhận bằng kiểm thử tự động; thực nghiệm chưa được chạy lại nên các số liệu trên chưa phản ánh chúng.

\clearpage

% ---------- Mục lục ----------
\tableofcontents
\clearpage
\pagenumbering{arabic}

\section{Bài toán đặt ra}

Khi chọn mua sách, người đọc thường phải tổng hợp nhiều thông tin cùng lúc: sách có phù hợp nhu cầu không, tác giả là ai, giá bao nhiêu, người mua trước nhận xét thế nào và có điểm nào cần lưu ý. Các thông tin này thường nằm ở những nơi khác nhau. Trang giới thiệu sản phẩm cung cấp thông số sách, đánh giá người mua phản ánh trải nghiệm, còn nguồn thư mục hoặc nhà xuất bản giúp đối chiếu thông tin về tác phẩm.

Một câu hỏi đơn giản như ``Cuốn sách này giá bao nhiêu?'' chủ yếu cần tra cứu. Nhưng yêu cầu ``Tìm sách phù hợp ngân sách, so sánh nhận xét của người mua và giải thích vì sao nên chọn'' đòi hỏi nhiều bước hơn. Hệ thống phải tìm đúng sách, đọc đúng đánh giá của cuốn đó, cân nhắc các điều kiện và đưa ra kết luận có căn cứ. Nếu ghép nhầm đánh giá giữa hai cuốn sách hoặc tự bổ sung thông tin chưa có nguồn, câu trả lời có thể trôi chảy nhưng vẫn sai.

Đề tài xây dựng một trợ lý tư vấn sách tiếng Việt theo hướng chia công việc cho nhiều tác tử chuyên trách. Trong báo cáo này, \textbf{tác tử} được hiểu là một thành phần đảm nhiệm một nhóm nhiệm vụ nhất định, chẳng hạn tìm thông tin sách hoặc phân tích đánh giá. Các tác tử được điều phối để cùng giải quyết yêu cầu; không nên hiểu rằng mỗi tác tử là một trí tuệ tự chủ hoàn toàn độc lập.

Vấn đề nghiên cứu không chỉ là làm cho trợ lý trả lời được. Cần làm rõ câu trả lời dựa vào đâu, các phần việc được phối hợp thế nào và việc dùng nhiều tác tử có đem lại lợi ích đủ lớn so với cách dùng một tác tử hay không.

\section{Mục tiêu và phạm vi}

\subsection{Mục tiêu của hệ thống}
Hệ thống hướng tới bốn mục tiêu chính:
\begin{enumerate}
\item \textbf{Hỗ trợ tìm và so sánh sách bằng tiếng Việt.} Người dùng có thể nêu nhu cầu, ngân sách hoặc những thông tin muốn đối chiếu bằng cách hỏi tự nhiên.
\item \textbf{Kết hợp nhiều góc nhìn.} Thông tin sách, đánh giá người mua và nguồn kiến thức được xem xét theo từng nhiệm vụ rồi tổng hợp thành câu trả lời.
\item \textbf{Giữ câu trả lời có căn cứ.} Thông tin được trình bày phải gắn với dữ liệu đã có; thiếu căn cứ thì cần nói rõ hoặc hỏi lại.
\item \textbf{Đánh giá lợi ích và chi phí của nhiều tác tử.} So sánh mức đáp ứng câu hỏi, thời gian trả lời và lượng tài nguyên sử dụng với cách dùng một tác tử.
\end{enumerate}

\subsection{Phạm vi sử dụng dữ liệu}
Hệ thống dùng dữ liệu sách và đánh giá người mua được thu thập trước đó, cùng một tập nguồn kiến thức được chọn lọc. Đây là \textbf{dữ liệu lịch sử}, không phải thông tin được cập nhật trực tiếp từ sàn thương mại điện tử. Vì vậy, giá trong câu trả lời phải được hiểu là giá ghi nhận trong dữ liệu, không mặc nhiên là giá đang bán ở thời điểm người dùng hỏi. \src{1}\src{2}

Hệ thống cũng không xác nhận được tình trạng tồn kho thực tế, doanh số hiện tại hoặc đánh giá nào là giả. Những nhận xét như ``có phàn nàn về đóng gói'' chỉ phản ánh nội dung các đánh giá đã được lấy mẫu. Chúng không đủ để kết luận toàn bộ người mua có cùng trải nghiệm.

Phần mở rộng của hệ thống có mô phỏng thao tác dành cho người mua và người bán, như chuẩn bị giỏ hàng hoặc đề xuất thay đổi thông tin bán hàng. Các thao tác này diễn ra trong môi trường thử nghiệm, không phải giao dịch thật. Trọng tâm báo cáo là cách tư vấn, phối hợp tác tử và kết quả thực nghiệm.

\section{Mô hình hoạt động của hệ thống}

\subsection{Ý tưởng chia công việc cho các tác tử}
Thay vì giao mọi việc cho một thành phần xử lý chung, hệ thống phân chia yêu cầu thành những phần có trách nhiệm rõ ràng. Một bộ phận điều phối nhận câu hỏi, xác định cần những thông tin nào và chuyển từng phần việc cho nhóm chuyên trách. Kết quả được kiểm tra trước khi ghép thành câu trả lời cuối. \src{1}

Ví dụ, với câu hỏi ``Nên chọn cuốn A hay cuốn B nếu tôi quan tâm đến nội dung và chất lượng đóng gói?'', việc so sánh thông số sách khác với việc đọc đánh giá đóng gói. Tách hai phần giúp hệ thống biết mỗi kết luận đến từ nguồn nào, đồng thời giảm nguy cơ lấy ý kiến về cuốn A để nhận xét cuốn B.

\begin{table}[htbp]\small
\caption{Các nhóm nhiệm vụ trong hệ thống.}
\begin{tabularx}{\linewidth}{P{33mm}YY}\toprule
Nhóm nhiệm vụ & Trách nhiệm & Ví dụ thông tin cung cấp\\\midrule
Thông tin sách & Tìm, lọc và so sánh các cuốn sách & Tên sách, tác giả, giá trong dữ liệu, số trang, điểm đánh giá\\
Đánh giá người mua & Tổng hợp ý kiến từ những đánh giá được lưu & Nội dung được khen, trải nghiệm đọc, nhận xét giao hàng\\
Phàn nàn và độ tin cậy & Nhận diện vấn đề và đặc điểm văn bản đánh giá & Đóng gói móp, sách hư hỏng, nội dung đánh giá quá ít thông tin\\
Tổng quan danh mục & Tổng hợp số liệu trong tập dữ liệu & So sánh theo tác giả, nhóm sách hoặc khoảng giá\\
Kiến thức về tác phẩm & Đối chiếu với nguồn kiến thức được chọn lọc & Tác giả gốc, tên tác phẩm gốc, thông tin được nguồn xác nhận\\\bottomrule
\end{tabularx}
\end{table}

Nhóm phàn nàn và độ tin cậy chỉ cung cấp \textbf{dấu hiệu tham khảo}. Một nhận xét ngắn chưa chắc là nhận xét không đáng tin, và việc phát hiện từ ngữ tiêu cực không đồng nghĩa đã phát hiện gian lận. Hệ thống phải giữ rõ giới hạn này khi diễn giải kết quả.

\subsection{Vai trò của bộ phận điều phối}
Bộ phận điều phối trả lời ba câu hỏi: người dùng cần biết gì, cần giao việc cho nhóm nào và các kết quả đã đủ để trả lời chưa. Không phải câu hỏi nào cũng cần kích hoạt mọi tác tử. Câu hỏi chỉ về giá có thể chỉ cần thông tin sách; câu hỏi về so sánh và trải nghiệm người mua mới cần thêm các góc nhìn khác.

Nếu chưa xác định được cuốn sách hoặc câu hỏi có điều kiện mâu thuẫn, hệ thống cần làm rõ trước. Chẳng hạn, yêu cầu ``giá trên 300.000 đồng nhưng dưới 200.000 đồng'' không nên được xử lý bằng cách tự bỏ một trong hai điều kiện. Việc hỏi lại đúng lúc là một phần của chất lượng tư vấn.

\fig{tong-quan}{Mô hình xử lý ở mức chức năng: hiểu yêu cầu, phối hợp chuyên gia, lấy bằng chứng rồi kiểm tra câu trả lời.}{fig:overview}

\subsection{Vai trò của bước tổng hợp}
Kết quả của từng nhóm chuyên trách chỉ là thông tin từng phần. Bước tổng hợp cần biến chúng thành một câu trả lời thống nhất, giải thích được vì sao một cuốn sách phù hợp hơn, đồng thời giữ lại những hạn chế của dữ liệu.

Ví dụ, sách A có điểm đánh giá cao nhưng trong mẫu có phàn nàn về đóng gói. Kết luận phù hợp không phải ``sách A chắc chắn không tốt'', mà là ``sách A đáp ứng điều kiện đã nêu; tuy nhiên một số đánh giá trong dữ liệu phản ánh vấn đề đóng gói''. Cách diễn đạt này giúp người đọc hiểu cả điểm mạnh lẫn mức chắc chắn của nhận định.

\section{Dữ liệu được sử dụng}

\subsection{Thông tin sách và đánh giá người mua}
Dữ liệu chính được lấy từ bộ dữ liệu Tiki Books đã công bố. Hệ thống sử dụng hai tập được chuẩn bị cho mục đích kiểm thử và đánh giá, thay vì lấy toàn bộ dữ liệu vào mọi lần chạy. \src{2}

\begin{table}[htbp]\small
\caption{Quy mô hai tập dữ liệu đã lưu trong dự án.}
\begin{tabularx}{\linewidth}{YrrY}\toprule
Tập dữ liệu & Bản ghi sản phẩm & Đánh giá & Mục đích\\\midrule
Tập kiểm thử & 24 & 115 & Kiểm tra chức năng trên một tập nhỏ\\
Tập đánh giá & 200 & 1.773 & Làm nền cho việc đánh giá hệ thống\\\bottomrule
\end{tabularx}
\end{table}

Đánh giá được lấy mẫu theo từng sản phẩm, không phải toàn bộ ý kiến đã từng xuất hiện trên sàn. Một sản phẩm có thể có nhiều lượt đánh giá trong thông tin tổng quan, nhưng hệ thống chỉ đọc được những đánh giá thực sự có trong tập dữ liệu đã lưu. Hai con số này cần được phân biệt khi trả lời.

Trước khi sử dụng, dữ liệu được làm sạch, loại trùng, chuẩn hóa và che những thông tin như số điện thoại, thư điện tử hoặc liên kết trong nội dung đánh giá. Khi kiểm đếm lại các tập đã lưu, số bản ghi khớp với mô tả, các mã sản phẩm và đánh giá không trùng nhau, mọi đánh giá đều gắn với một sản phẩm có trong tập tương ứng. Điều đó hỗ trợ độ nhất quán của dữ liệu, nhưng chưa chứng minh mọi nội dung đều chính xác về mặt thực tế.

\subsection{Nguồn kiến thức về tác phẩm}
Bên cạnh thông tin bán hàng, hệ thống có một tập nguồn được chọn lọc để đối chiếu kiến thức về tác phẩm. Nguồn có thể đến từ tác giả, nhà xuất bản hoặc thư viện. Tập hiện có 20 nguồn, hỗ trợ đối chiếu ở cấp tác phẩm cho 20 trong 200 bản ghi của danh mục. \src{3}

Cần phân biệt \textbf{tác phẩm} và \textbf{ấn bản}. Một nguồn có thể xác nhận ai viết tác phẩm, nhưng chưa đủ xác nhận cuốn đang bán thuộc lần tái bản nào, có bao nhiêu trang hoặc do ai dịch. Nếu gộp các thông tin này, hệ thống dễ tạo ra một câu trả lời có vẻ đầy đủ nhưng sai đối tượng.

Ví dụ, tài liệu đối chiếu của dự án ghi nhận một nguồn thư mục về sách Morisaki có số trang khác với dữ liệu bán hàng. Cách xử lý đúng là nhận biết hai nguồn chưa xác nhận cùng một ấn bản, không tự chọn một con số rồi coi đó là sự thật chắc chắn. \src{3}

\begin{table}[htbp]\small
\caption{Mức độ đối chiếu của nguồn kiến thức với danh mục.}
\begin{tabularx}{\linewidth}{YrY}\toprule
Tình trạng & Số bản ghi & Cách hiểu\\\midrule
Đối chiếu được ở cấp tác phẩm & 20 & Có nguồn hỗ trợ nhận diện tác phẩm; chưa xác nhận đầy đủ ấn bản\\
Còn mơ hồ & 17 & Có thể là bộ sách hoặc nhiều tác phẩm ghép trong một sản phẩm\\
Chưa có đối chiếu trong tập nguồn & 163 & Chưa được bổ sung nguồn phù hợp trong phạm vi tuyển chọn hiện tại\\\bottomrule
\end{tabularx}
\end{table}

Như vậy, phạm vi kiến thức đối chiếu được còn hẹp. ``Chưa có nguồn trong hệ thống'' không có nghĩa ngoài thực tế không có nguồn, mà chỉ có nghĩa hệ thống chưa đủ dữ liệu để trả lời có căn cứ. Ngoài ra, dự án ghi nhận một bản ghi bán lẻ không thuộc sách trong danh mục. Đây là điểm cần tiếp tục làm sạch khi mở rộng dữ liệu. \src{3}

\subsection{Dữ liệu có thể và không thể chứng minh điều gì}
\begin{table}[htbp]\small
\caption{Giới hạn diễn giải theo từng loại thông tin.}
\begin{tabularx}{\linewidth}{P{33mm}YY}\toprule
Thông tin & Có thể sử dụng để & Không đủ để khẳng định\\\midrule
Giá trong dữ liệu & So sánh mức giá đã được ghi nhận & Giá đang bán hôm nay\\
Đánh giá lấy mẫu & Mô tả ý kiến xuất hiện trong mẫu & Ý kiến của toàn bộ người mua\\
Chỉ báo phổ biến & Tham khảo mức phổ biến của dữ liệu lịch sử & Doanh số hiện tại đã được xác minh\\
Nguồn về tác phẩm & Kiểm tra nội dung mà nguồn thực sự hỗ trợ & Mọi chi tiết của ấn bản đang được bán\\\bottomrule
\end{tabularx}
\end{table}

\section{Quy trình xử lý một yêu cầu tư vấn}

\subsection{Bước 1: hiểu yêu cầu và điều kiện}
Hệ thống xác định người dùng muốn tìm sách, so sánh, xem đánh giá hay hỏi kiến thức. Đồng thời, các điều kiện như ngân sách, số lượng sách và tiêu chí quan tâm cần được giữ lại. Nếu người dùng hỏi nhiều ý, hệ thống phải nhận diện đủ các ý thay vì chỉ trả lời phần dễ nhất.

Ví dụ: ``Tìm hai cuốn sách dưới 200.000 đồng và cho biết người mua có phàn nàn về đóng gói không.'' Yêu cầu có ba phần: chọn đúng số lượng, đáp ứng điều kiện giá và kiểm tra nhận xét đóng gói. Một danh sách tên sách đơn thuần chưa hoàn thành yêu cầu.

\subsection{Bước 2: tìm và xác định đúng sách}
Nhóm thông tin sách tìm các ứng viên phù hợp, giữ tên đầy đủ và những thông tin cần thiết để phân biệt chúng. Đây là bước nền tảng: nếu chọn nhầm cuốn sách hoặc nhầm ấn bản, các phân tích phía sau dù đúng với dữ liệu đã lấy vẫn có thể trả lời sai câu hỏi.

Khi cần xếp thứ tự ứng viên, hệ thống tham khảo điểm đánh giá, mức phổ biến lịch sử và giá tương đối. Những tiêu chí này giúp tạo thứ tự lựa chọn rõ ràng, nhưng không phải một phép đo đầy đủ về giá trị văn học hay mức phù hợp với mọi người đọc. Giá rẻ hơn hoặc nhiều lượt đánh giá hơn không tự động có nghĩa là sách hay hơn.

\subsection{Bước 3: bổ sung đánh giá và bằng chứng}
Sau khi xác định sách, các nhóm chuyên trách xem xét đánh giá người mua, dấu hiệu phàn nàn và nguồn kiến thức nếu câu hỏi cần. Mọi nhận xét phải đi cùng đúng cuốn sách đã chọn. Những ý kiến về đóng gói hay giao hàng cũng cần tách khỏi nhận xét về chất lượng nội dung tác phẩm.

\fig{phoi-hop}{Ví dụ phối hợp khi tư vấn lựa chọn sách. Các nhánh thể hiện các phần việc cần kết hợp, không hàm ý chúng luôn được thực hiện đồng thời.}{fig:cooperation}

\subsection{Bước 4: tổng hợp và kiểm tra}
Hệ thống ghép các kết quả thành một câu trả lời thống nhất, kiểm tra những ý đã có bằng chứng và những phần chưa được hỗ trợ. Một nguồn liên quan đến cuốn sách chưa đủ để chứng minh mọi điều được viết về cuốn sách đó; nguồn phải thực sự hỗ trợ nội dung đang được nêu.

Nếu một phần thông tin chưa có, hệ thống có thể trả lời phần đã biết và chỉ rõ phần còn thiếu. Nếu thông tin thiếu là điều kiện quyết định, chẳng hạn chưa biết người dùng đang hỏi cuốn nào, hệ thống nên hỏi lại. Không nên điền một thông tin có vẻ hợp lý chỉ để câu trả lời trông hoàn chỉnh.

\subsection{Bước 5: trình bày kết quả cho người dùng}
Một câu trả lời phù hợp nên đi theo trình tự: trả lời trực tiếp, đưa thông tin so sánh cần thiết, giải thích căn cứ và nêu giới hạn còn lại. Khi so sánh nhiều cuốn, cần trình bày song song để người đọc thấy điểm khác biệt, không trộn tất cả nhận xét thành một đoạn chung khó đối chiếu.

\section{Cơ chế giữ chất lượng câu trả lời}

\subsection{Giới hạn phạm vi công việc}
Hệ thống giới hạn số ứng viên và số lần bổ sung thông tin trong một yêu cầu. Mục đích là tránh xử lý kéo dài không kiểm soát và hạn chế việc càng suy luận càng đi xa dữ liệu. Với yêu cầu quá rộng, cách phù hợp là chia nhỏ hoặc hỏi thêm tiêu chí, thay vì cố liệt kê và đánh giá mọi sách cùng lúc.

\subsection{Giữ đúng nguồn của từng kết luận}
Thông tin danh mục trả lời các câu hỏi về sản phẩm được lưu; đánh giá phản ánh ý kiến trong mẫu; nguồn kiến thức hỗ trợ những thông tin đã được đối chiếu. Bước tổng hợp cần giữ ranh giới giữa ba loại này. Chẳng hạn, ``người mua khen nội dung'' là một nhận xét trong đánh giá, không phải kết luận khách quan rằng tác phẩm phù hợp với tất cả người đọc.

\subsection{Không tự thực hiện thay đổi ngoài thẩm quyền}
Phần tư vấn không được tự sửa dữ liệu hoặc biến một lời đề nghị trong hội thoại thành giao dịch đã hoàn tất. Với các thao tác mô phỏng có làm thay đổi thông tin, hệ thống cần tạo đề xuất và có bước xác nhận phù hợp. Việc tư vấn và việc thực hiện hành động phải được phân biệt để người dùng biết điều gì mới là gợi ý, điều gì đã được thực hiện.

\subsection{Sửa các điểm làm mất thông tin}
Sau thực nghiệm, nhóm rà soát đường đi của thông tin từ câu hỏi đến câu trả lời và tái hiện được sáu chỗ hệ thống có thể làm rơi dữ kiện. Mỗi chỗ được sửa riêng và áp dụng chung cho cả cách một tác tử lẫn nhiều tác tử. \src{7}
\begin{enumerate}
\item Khi mô hình viết lại câu tìm kiếm, tên sách hoặc mã người dùng nêu có thể bị mất. Nay câu hỏi gốc luôn được giữ cùng các cách viết lại.
\item Nguồn chỉ nhắc tên sách ở tiêu đề có thể không được tìm thấy. Nay tiêu đề được dùng để tìm nguồn, nhưng không được coi là bằng chứng cho nội dung.
\item Cùng một chữ tiếng Việt có dấu có thể được mã hóa theo hai cách khác nhau, khiến kết quả tìm kiếm khác nhau. Nay văn bản được chuẩn hóa trước khi so khớp.
\item Hai nguồn khác nhau có cùng đoạn văn từng bị gộp làm một, mất ghi nhận nguồn. Nay chỉ gộp các đoạn trùng trong cùng một nguồn.
\item Đoạn trích bị cắt ở 350 ký tự có thể mất chi tiết ở cuối. Nay đoạn trích được giữ tới 1.000 ký tự khi còn chỗ, trong cùng giới hạn đầu vào.
\item Câu trả lời từng được coi là đủ chỉ vì có một thông tin đúng. Nay mỗi cuốn được gợi ý phải có tên sách, kèm giá, điểm đánh giá hoặc số trang nếu người dùng lọc theo các điều kiện đó.
\end{enumerate}
Các sửa đổi được kiểm tra bằng tám tình huống kiểm thử dựng riêng cho từng lỗi. Trên mã nguồn trước khi sửa, sáu tình huống thất bại; trên mã nguồn hiện tại, cả tám tình huống đều qua. Đây là bằng chứng các lỗi cụ thể đã được khắc phục, \textbf{không phải điểm chất lượng của hệ thống}. Thực nghiệm ở mục sau được chạy trước đợt sửa này nên chưa phản ánh tác động của chúng.

\subsection{Kiểm tra tự động vẫn có giới hạn}
Hệ thống có bước kiểm tra câu trả lời, nhưng bước kiểm tra cũng có thể đánh giá sai. Vì thế, có cơ chế kiểm tra không đồng nghĩa bảo đảm tuyệt đối không có thông tin bịa hoặc nhầm. Độ tin cậy cuối cùng cần được đánh giá bằng các câu hỏi cụ thể và đối chiếu với nguồn, thay vì suy ra từ tên gọi của một thành phần.

\section{Thiết kế thực nghiệm}

\subsection{Mục đích so sánh}
Thực nghiệm so sánh hai cách trả lời cùng một bộ câu hỏi:
\begin{itemize}
\item \textbf{Một tác tử:} nhận câu hỏi và thông tin liên quan, sau đó trực tiếp tạo câu trả lời.
\item \textbf{Nhiều tác tử:} phân chia công việc cho các chuyên gia cần thiết, rồi tổng hợp kết quả thành một câu trả lời chung.
\end{itemize}
Cả hai cách dùng cùng một mô hình ngôn ngữ nền. Do đó, khác biệt cần quan sát chủ yếu nằm ở cách tổ chức xử lý và cách hướng dẫn từng phần việc, không phải việc một bên được dùng mô hình nền mạnh hơn. \src{4}\src{5}

Thực nghiệm này là một phép so sánh riêng, có điều kiện được kiểm soát. Nó chưa đo toàn bộ quá trình sử dụng ứng dụng trong thực tế. Giữ rõ phạm vi này giúp đánh giá đúng giá trị của kết quả.

\subsection{Bộ câu hỏi}
Bộ thử nghiệm gồm 24 câu hỏi được chia thành bốn nhóm. Mỗi câu có một hoặc nhiều dữ kiện bắt buộc mà câu trả lời cần đề cập đúng. Tổng cộng có 43 dữ kiện bắt buộc cho một lượt đi qua cả bộ câu hỏi. \src{4}

\begin{table}[htbp]\small
\caption{Cấu trúc bộ câu hỏi thực nghiệm.}
\begin{tabularx}{\linewidth}{P{12mm}Yrr}\toprule
Nhóm & Nội dung chính & Câu hỏi & Dữ kiện\\\midrule
1 & Tra cứu chủ yếu thuộc một chuyên gia & 6 & 14\\
2 & Kết hợp thông tin của nhiều chuyên gia & 8 & 17\\
3 & Tư vấn lựa chọn theo nhiều góc nhìn & 6 & 8\\
4 & Nhận biết yêu cầu không phù hợp hoặc cần từ chối & 4 & 4\\\midrule
Tổng & & 24 & 43\\\bottomrule
\end{tabularx}
\end{table}

Mỗi câu được chạy ba lần với mỗi cách xử lý. Tổng số câu trả lời được lưu là $24\times2\times3=144$, tức 72 câu trả lời cho mỗi bên. Khi chấm các dữ kiện bắt buộc, mỗi bên có $43\times3=129$ lượt đánh giá dữ kiện.

\textbf{144 lượt chạy không phải 144 bài toán khác nhau.} Thực chất vẫn là 24 câu hỏi được lặp lại. Sự phân biệt này quan trọng khi nhận xét độ đa dạng và mức chắc chắn của kết quả.

\subsection{Điều kiện cung cấp thông tin}
Trước khi trả lời, cả hai cách đều được cấp sẵn những thông tin liên quan đến câu hỏi. Việc chọn thông tin dựa vào nguồn của đáp án chuẩn và một số bổ sung đã được người thiết kế xác định. Điều này giúp hai bên có cùng nền dữ liệu, nhưng cũng làm bước tìm kiếm ban đầu thuận lợi hơn so với sử dụng thực tế. \src{5}

Vì vậy, phép thử trả lời câu hỏi: \textit{khi đã có những thông tin liên quan, cách tổ chức nào đáp ứng yêu cầu tốt hơn?} Nó chưa trả lời đầy đủ câu hỏi: \textit{cách nào tự tìm đúng thông tin từ toàn bộ kho dữ liệu tốt hơn?}

\fig{so-sanh-thuc-nghiem}{Hai cách xử lý được so sánh trên thông tin liên quan đã cấp sẵn. Bộ chấm đối chiếu câu trả lời với các dữ kiện bắt buộc.}{fig:experiment}

\subsection{Cách chấm điểm và các tiêu chí quan sát}
Một mô hình chấm tự động đối chiếu câu trả lời với từng dữ kiện bắt buộc. Điểm của một câu là tỷ lệ dữ kiện được chấp nhận trong câu đó. Điểm trung bình toàn bộ câu hỏi được tính bằng cách lấy trung bình qua các lần lặp rồi qua 24 câu, để mỗi câu hỏi có trọng số như nhau. \src{5}

Báo cáo đồng thời nêu tỷ lệ dữ kiện được chấp nhận khi gộp toàn bộ 129 lượt đánh giá. Hai cách tính này có thể khác nhau vì mỗi câu hỏi có số dữ kiện khác nhau. Một câu có bốn dữ kiện sẽ có ảnh hưởng lớn hơn trong cách gộp theo dữ kiện, nhưng vẫn chỉ có trọng số như các câu khác trong cách tính trung bình theo câu hỏi.

Ngoài điểm nội dung, thực nghiệm ghi thời gian tạo câu trả lời và lượng văn bản mô hình xử lý. Lượng văn bản được ghi bằng \textit{token}, là đơn vị dùng để biểu diễn độ dài đầu vào và đầu ra của mô hình. Token không tương đương chính xác với số từ, nhưng có ích để so sánh mức sử dụng tài nguyên giữa hai cách.

\section{Kết quả thực nghiệm}

\subsection{Kết quả tổng hợp}
Các con số dưới đây được tính lại từ những câu trả lời, điểm chấm và thông tin đo lường đã lưu của thực nghiệm; không phải một đợt chạy mô hình mới. \src{6}

\begin{table}[htbp]\small
\caption{So sánh hai cách xử lý trên bộ thực nghiệm.}
\begin{tabularx}{\linewidth}{Yrr}\toprule
Tiêu chí & Một tác tử & Nhiều tác tử\\\midrule
Số câu trả lời được lưu & 72 & 72\\
Điểm trung bình theo câu hỏi & 98,96\% & 100,00\%\\
Lượt dữ kiện được chấp nhận & 126/129 & 129/129\\
Tỷ lệ chấp nhận theo dữ kiện & 97,67\% & 100,00\%\\
Thời gian trả lời trung bình & 1,43 giây & 6,75 giây\\
Token trung bình mỗi câu trả lời & 1.332 & 3.706\\\bottomrule
\end{tabularx}
\end{table}

Nhiều tác tử đạt điểm trung bình cao hơn khoảng \textbf{1,04 điểm phần trăm}. Tuy nhiên, thời gian trả lời trung bình dài gấp khoảng \textbf{4,73 lần}, còn lượng token lớn gấp khoảng \textbf{2,78 lần}. Như vậy, phần cải thiện nội dung quan sát được đi kèm mức tăng đáng kể về thời gian và tài nguyên.

Điểm 100\% ở đây có nghĩa tất cả dữ kiện bắt buộc đều được bộ chấm chấp nhận. Nó \textbf{không có nghĩa toàn bộ hệ thống chính xác tuyệt đối}, vì bộ chấm chưa đối chiếu mọi phát biểu bổ sung trong câu trả lời với toàn bộ dữ liệu nguồn.

\subsection{Kết quả theo nhóm câu hỏi}
\begin{table}[htbp]\small
\caption{Điểm và thời gian trung bình theo nhóm câu hỏi.}
\begin{tabularx}{\linewidth}{Yrrrr}\toprule
 & \multicolumn{2}{c}{Điểm trung bình (\%)} & \multicolumn{2}{c}{Thời gian (giây)}\\
Nhóm & Một tác tử & Nhiều tác tử & Một tác tử & Nhiều tác tử\\\midrule
1. Tra cứu đơn giản & 95,83 & 100,00 & 1,35 & 5,29\\
2. Kết hợp nhiều nguồn thông tin & 100,00 & 100,00 & 1,50 & 8,82\\
3. Tư vấn nhiều góc nhìn & 100,00 & 100,00 & 1,52 & 8,40\\
4. Yêu cầu không phù hợp & 100,00 & 100,00 & 1,26 & 2,32\\\bottomrule
\end{tabularx}
\end{table}

Ở nhóm 2 và 3, cả hai cách cùng đạt điểm tối đa theo dữ kiện bắt buộc. Do đó, bộ thử nghiệm này chưa cho thấy nhiều tác tử tạo được chênh lệch điểm ở chính các câu hỏi được thiết kế theo hướng kết hợp chuyên gia. Việc gắn một câu hỏi vào nhóm phức tạp không tự chứng minh rằng chỉ nhiều tác tử mới trả lời được.

\subsection{Trường hợp tạo ra chênh lệch}
Toàn bộ khác biệt điểm nằm ở một câu hỏi về \textit{Sapiens}. Tên sách đầy đủ trong dữ liệu có hậu tố \textit{(Tái Bản 2022)}. Cách một tác tử bỏ hậu tố này trong cả ba lần trả lời; cách nhiều tác tử giữ được tên đầy đủ. Vì vậy, một tác tử đạt ba trên bốn dữ kiện của câu hỏi đó, còn nhiều tác tử đạt đủ bốn. \src{6}

Xét trên 24 câu hỏi, nhiều tác tử có điểm cao hơn ở một câu và hòa ở 23 câu còn lại. Kết quả này cho thấy tầm quan trọng của việc giữ đúng thông tin ấn bản, nhưng chưa đủ để khẳng định nhiều tác tử tốt hơn trên mọi dạng câu hỏi.

Cần lưu ý rằng chuyên gia thông tin sách được hướng dẫn nhấn mạnh tên đầy đủ và hậu tố ấn bản, trong khi cách một tác tử được yêu cầu trả lời ngắn gọn hơn. Do đó, ưu thế ở câu này có thể đến từ cách hướng dẫn trình bày chi tiết, chứ chưa chắc chỉ do việc chia thành nhiều tác tử. Một phép đối chứng hợp lý là yêu cầu cách một tác tử cũng bảo toàn tên ấn bản rồi thử trên một bộ câu hỏi mới. \src{5}

\subsection{Mức chắc chắn của kết luận}
Khoảng ước lượng cho chênh lệch điểm trung bình còn bao gồm cả khả năng không có cải thiện. Kết hợp với việc chỉ một trong 24 câu tạo ra khác biệt, dữ liệu hiện có \textbf{chưa đủ để kết luận nhiều tác tử có ưu thế tổng quát rõ ràng}. Điều này cũng không chứng minh hai cách hoàn toàn tương đương; cần thêm câu hỏi và một thiết kế đánh giá đủ sức phân biệt chúng.

Việc cả hai bên đạt điểm rất cao còn tạo ra một giới hạn đo lường: khi một tác tử đã gần điểm tối đa, nhiều tác tử dù làm tốt cũng khó tạo thêm nhiều chênh lệch. Muốn đánh giá sâu hơn, cần các câu hỏi có khó khăn thực chất về chọn đúng sách, đối chiếu nguồn hoặc xử lý thông tin thiếu và mâu thuẫn.

\subsection{Đối chiếu với một bộ câu hỏi khó hơn}
Trước bộ 24 câu ở trên, dự án đã chạy một đợt đánh giá lớn hơn gồm 60 hội thoại thuộc 27 nhóm tác phẩm, mỗi hội thoại chạy ba lần với bốn cấu hình, tổng cộng 720 lượt quan sát. Cả 720 lượt đều được chấm và gói kết quả gốc qua kiểm tra toàn vẹn tệp.

Trên bộ này, tỷ lệ hoàn thành tác vụ trung bình thấp hơn nhiều: khoảng 45,2\% (một tác tử), 47,2\% (nhiều tác tử, kế hoạch cố định) và 41,4\% (nhiều tác tử, kế hoạch thích nghi). So sánh chính đã đăng ký trước, giữa một tác tử và nhiều tác tử thích nghi, cho chênh lệch $-3{,}8$ điểm phần trăm với khoảng tin cậy 95\% từ $-8{,}7$ đến $+1{,}0$; tức không có bằng chứng nhiều tác tử tốt hơn, và ước lượng điểm còn nghiêng nhẹ về một tác tử.

Phần đóng góp rõ nhất đến từ việc truy xuất có căn cứ. Khi tắt truy xuất trong cùng cấu hình nhiều tác tử thích nghi, tỷ lệ hoàn thành giảm khoảng 36,1 điểm phần trăm (khoảng tin cậy 95\% từ $-51{,}2$ đến $-21{,}6$) và khả năng từ chối đúng lúc giảm khoảng 74,4 điểm phần trăm. Cả bốn cấu hình đều không ghi nhận vi phạm phân quyền.

Cần đọc các số này với các giới hạn: bộ chấm chưa được hiệu chuẩn bằng người thật, khả năng tìm đúng tài liệu (document recall) chưa được đo, và bộ 60 hội thoại này đã được chạy lặp nhiều lần trong quá trình phát triển, hệ thống được sửa giữa các lần, nên không thể xem đây là phép xác nhận độc lập cho phiên bản hiện tại.

\subsection{Lần chạy xác nhận trên câu hỏi mới}
Để có một phép thử chưa từng dùng để chỉnh hệ thống, dự án soạn 60 hội thoại mới trên cùng 13 sản phẩm và 26 nhóm tác phẩm, đóng băng bằng mã băm, rồi chạy đủ bốn cấu hình, ba lần lặp, tổng cộng 720 lượt. Giả thuyết và quy tắc phân tích được ghi lại và lưu vào kho mã \textit{trước khi} soạn câu hỏi và chạy hệ thống: giả thuyết chính là nhiều tác tử kế hoạch cố định hoàn thành tác vụ tốt hơn một tác tử ở loại câu nhiều điều kiện; chỉ kết luận ``ủng hộ'' khi khoảng tin cậy 95\% (bootstrap theo nhóm tác phẩm, 10.000 mẫu) không chứa 0. \src{10}

\begin{table}[htbp]\small
\caption{Kết quả lần chạy xác nhận: chênh lệch tỷ lệ hoàn thành tác vụ so với một tác tử (điểm phần trăm).}
\begin{tabularx}{\linewidth}{Yrrr}\toprule
So sánh & Chênh lệch & Khoảng tin cậy 95\% & Thắng/Hòa/Thua\\\midrule
Kế hoạch cố định, câu nhiều điều kiện (giả thuyết chính) & $+2{,}78$ & $[0{,}00;\ +8{,}33]$ & 1/11/0\\
Kế hoạch cố định, cả 60 hội thoại & $+1{,}07$ & $[0{,}00;\ +2{,}78]$ & 2/24/0\\
Kế hoạch thích nghi, cả 60 hội thoại & $+1{,}07$ & $[0{,}00;\ +2{,}78]$ & 2/24/0\\
Thích nghi không truy xuất, cả 60 hội thoại & $-16{,}24$ & $[-23{,}18;\ -9{,}62]$ & 0/13/13\\\bottomrule
\end{tabularx}
\end{table}

\textbf{Không giả thuyết nào được ủng hộ}: cả ba khoảng tin cậy của so sánh nhiều tác tử với một tác tử đều chạm 0. Điều nói được là hai cấu hình nhiều tác tử không thua một tác tử ở nhóm tác phẩm nào và cao hơn khoảng 1,07 điểm phần trăm, nhưng chênh lệch chỉ đến từ 2 trong 26 nhóm và chưa có ý nghĩa thống kê. Tắt truy xuất tiếp tục làm giảm rõ tỷ lệ hoàn thành, nhất quán với đợt đánh giá trước.

Tỷ lệ hoàn thành tuyệt đối trên bộ mới thấp: 19,4\% (một tác tử) và 20,6\% (hai cấu hình nhiều tác tử). Bốn trong sáu loại câu cho 0\% ở mọi cấu hình, và loại câu nhiều điều kiện giảm từ khoảng 94\% ở đợt trước xuống khoảng 6\%. Ở các loại này, cả một tác tử lẫn nhiều tác tử thường trả lời rằng chưa tìm thấy đủ ứng viên và đề nghị người dùng nêu rõ hơn. Câu hỏi của đợt trước theo một khuôn cố định, còn câu hỏi mới dùng đúng tên sách nhưng diễn đạt tự nhiên hơn; bước lập kế hoạch và tìm ứng viên là phần dùng chung của hai cách nên cả hai gặp khó giống nhau. Cách hiểu hợp lý nhất, dù mới dựa trên một số ví dụ, là khả năng xử lý cách hỏi ngoài khuôn quen thuộc còn yếu, và điểm cao ở đợt trước một phần phản ánh việc hệ thống đã được chỉnh nhiều vòng trên chính khuôn câu hỏi đó. Hệ quả là bộ mới chỉ phân biệt được hai cách ở loại câu tra cứu nguồn kiến thức (91,7\% so với 94,4\% và 97,2\%).

Kết quả này có các giới hạn phải nêu: câu hỏi mới dùng lại cùng 13 sản phẩm và do một mô hình ngôn ngữ soạn, nên không phải bộ độc lập về dữ liệu; ngưỡng hiệu chuẩn bộ chấm đã được nới vì bộ chấm lệch nhãn tham chiếu tự động ở 3 trong 32 trường hợp thử; 719 trên 720 lượt được chấm, lượt còn lại bộ chấm không trả được kết quả hợp lệ và không thuộc các so sánh nhiều tác tử với một tác tử; bảng trên được tính trực tiếp từ nhật ký chấm theo công thức đã khai báo, không qua bước tổng hợp chính thức của quy trình. Câu hỏi và đáp án không được sửa sau khi xem kết quả.

\subsection{Chấm lại bằng cùng bộ chấm (thăm dò)}
Để đo độ ổn định của bộ chấm mô hình mà không cần người chấm, 720 câu trả lời của đợt lớn được chấm lại bằng đúng cấu hình cũ (cùng mô hình \texttt{gpt-6-luna}). Trên 715 câu có kết quả hợp lệ ở cả hai lần, độ đồng thuận là 97,3\% (hệ số Cohen $\kappa = 0{,}946$) với tiêu chí hoàn thành tác vụ và 96,9\% ($\kappa = 0{,}937$) với tiêu chí căn cứ của phát biểu. Chênh lệch giữa một tác tử và nhiều tác tử thích nghi là $-5{,}0$ điểm phần trăm ở lần chấm gốc và $-6{,}5$ điểm phần trăm ở lần chấm lại trên cùng 715 câu; hai lần chấm cùng chiều.

Kết quả này chỉ cho thấy so sánh ổn định qua hai lần chấm bằng cùng bộ chấm. Nó không chứng minh bộ chấm đúng: cùng một mô hình không phát hiện được thiên lệch chung, chế độ không suy luận có thể làm độ ồn bị đánh giá thấp, và bộ chấm vẫn chưa được người hiệu chuẩn. Kết luận chính thức vẫn là kết quả trên đủ 720 câu ở mục trước.

Với bộ 24 câu, chấm lại 144 câu trả lời bằng bộ chấm được cấp bản ghi nguồn cho điểm 88,5\% (một tác tử) và 96,5\% (nhiều tác tử), thay vì 99,0\% và 100\%. Kiểm tra thủ công cho thấy một phần mức giảm là do ngữ cảnh cấp cho bộ chấm chỉ gồm bản ghi của từng sự kiện, nên không dùng hai số này làm độ chính xác. Điều rút ra được là điểm bão hòa của bộ 24 câu phụ thuộc cách thiết kế bộ chấm và bộ này không đủ sức phân biệt hai kiến trúc. Chi tiết và giới hạn nằm trong \texttt{docs/rejudge\_llm\_judge\_results.md}.

\subsection{Kiểm chứng lại nguồn gốc số liệu}
Các số liệu trong mục này được tính lại độc lập từ kết quả đã lưu và khớp với bảng trên. Trong quá trình rà soát, nhóm cũng phát hiện một số tài liệu khác của dự án dùng số đo được viết sẵn trong mã chương trình, không tính từ kết quả chạy thật: báo cáo đánh đổi giữa các cấu hình, thí nghiệm bỏ bớt thành phần và một đợt chấm của người được mô phỏng. \src{8}

Các số liệu này \textbf{không được dùng trong báo cáo}. Những tài liệu liên quan đã được ghi chú đính chính; quy trình tạo báo cáo mới chỉ đọc kết quả gốc, kiểm tra tính toàn vẹn của tệp và dừng lại nếu thiếu dữ liệu, thay vì tự điền con số. Công cụ chấm của người cũng chỉ tổng hợp phiếu do người thật điền, không tự sinh nhãn.

Toàn bộ 1.456 kiểm thử tự động không gọi mô hình thật, gồm các kiểm thử chạy trên cơ sở dữ liệu PostgreSQL, đều qua trên mã nguồn hiện tại. Kết quả này cho thấy phần mềm vận hành đúng các quy tắc đã được kiểm thử; nó không thay thế một đợt đánh giá mới với mô hình thật.

\section{Nhận xét và thảo luận}

\subsection{Giá trị của việc chia trách nhiệm}
Lợi ích dễ nhận thấy của mô hình nhiều tác tử là tổ chức công việc rõ ràng. Hệ thống biết phần nào là tìm thông tin sách, phần nào là đọc đánh giá và phần nào là đối chiếu nguồn. Cách chia này giúp giải thích kết quả, tìm nguyên nhân khi sai và bổ sung năng lực chuyên môn theo từng phần.

Ví dụ, nếu câu trả lời chọn đúng sách nhưng nhận xét đóng gói sai, có thể xem lại phần đọc đánh giá thay vì sửa mọi bước. Nếu thông tin tác giả khác nhau giữa danh mục bán hàng và nguồn thư mục, hệ thống có cơ hội phát hiện sự khác biệt trước khi tổng hợp. Đây là giá trị về tổ chức và kiểm soát, dù chưa đồng nghĩa đã làm tăng điểm trên mọi bộ câu hỏi.

\subsection{Vì sao nhiều tác tử chậm hơn}
Cách một tác tử có thể đi trực tiếp từ thông tin được cấp đến câu trả lời. Cách nhiều tác tử còn phải phân chia yêu cầu, tạo kết quả từng phần và tổng hợp chúng. Mỗi công đoạn đều có thể phát sinh thêm thời gian và thêm văn bản cần xử lý.

Một số phần việc có thể được xử lý cùng lúc trong mô hình thử nghiệm, nhưng tổng thời gian vẫn còn những bước phải chờ nhau. Chẳng hạn, phải biết đã chọn cuốn sách nào thì mới có thể đọc đúng đánh giá của cuốn đó. Vì vậy, tăng số tác tử không tự động làm hệ thống nhanh hơn.

\subsection{Khi nào cần cân nhắc từng cách}
Với câu hỏi ngắn, dữ liệu liên quan đã có sẵn và ít yêu cầu đối chiếu, một tác tử có thể đáp ứng tốt với thời gian thấp hơn. Kết quả thực nghiệm hiện tại phù hợp với nhận định có điều kiện này.

Nhiều tác tử đáng cân nhắc khi một yêu cầu có nhiều loại nhiệm vụ cần kiểm soát riêng, cần đối chiếu các nguồn khác nhau hoặc cần giải thích rõ căn cứ của từng kết luận. Tuy nhiên, đây là lý do thiết kế và hướng ứng dụng; vẫn cần thực nghiệm phù hợp để xác nhận lợi ích chất lượng lớn hơn phần chi phí tăng thêm.

Không nên kết luận một cách tuyệt đối rằng ``câu hỏi khó thì bắt buộc nhiều tác tử''. Một tác tử cũng có thể được thiết kế để tra cứu từng bước và kiểm tra kết quả. Cần so sánh hai cách trên điều kiện tương đương trước khi quy toàn bộ khác biệt cho số lượng tác tử.

\section{Hạn chế của nghiên cứu hiện tại}

\subsection{Dữ liệu chưa đại diện cho sử dụng thực tế}
Thông tin sách là dữ liệu lịch sử, đánh giá người mua chỉ là mẫu và nguồn kiến thức mới phủ một phần nhỏ danh mục. Vì vậy, không thể dùng kết quả để khẳng định hệ thống hiểu toàn bộ thị trường sách hoặc luôn có thông tin cập nhật.

Phần chuẩn bị dữ liệu cũng cần tiếp tục làm sạch các bản ghi ngoài phạm vi sách và đối chiếu tốt hơn giữa tác phẩm với ấn bản. Chất lượng tư vấn phụ thuộc cả việc mô hình xử lý tốt lẫn việc dữ liệu đầu vào mô tả đúng đối tượng.

\subsection{Số câu hỏi còn ít}
Bộ thực nghiệm có 24 câu hỏi. Chạy mỗi câu ba lần giúp quan sát việc trả lời lặp lại, nhưng không tăng tương ứng số tình huống độc lập. Kết quả cần được kiểm tra trên các câu hỏi mới trước khi mở rộng thành nhận định chung.

Trong dữ liệu đã lưu, nội dung trả lời của ba lần lặp không hoàn toàn giống nhau. Do đó, việc đặt điều kiện sinh ít ngẫu nhiên không bảo đảm mô hình luôn tạo ra cùng một câu chữ. Cần đánh giá độ ổn định theo nội dung và dữ kiện, thay vì chỉ kỳ vọng văn bản lặp lại nguyên vẹn.

\subsection{Khả năng xử lý cách hỏi mới còn yếu}
Lần chạy xác nhận trên 60 hội thoại mới cho thấy cả một tác tử lẫn nhiều tác tử thường đề nghị người dùng nêu rõ hơn khi câu hỏi được diễn đạt tự nhiên thay vì theo khuôn của bộ câu hỏi cũ. Bước lập kế hoạch và tìm ứng viên là phần dùng chung, nên hạn chế này không phân biệt hai cách tổ chức; nó cho thấy điểm số trên các bộ câu hỏi quen thuộc không nên được hiểu là khả năng xử lý mọi cách hỏi. \src{10}

\subsection{Thông tin liên quan đã được cấp sẵn}
Hai cách trả lời đều nhận dữ liệu đã được lựa chọn dựa trên nguồn đáp án chuẩn. Đây là điều kiện hữu ích để so cách xử lý thông tin, nhưng chưa đánh giá đầy đủ khả năng tự tìm kiếm trong kho lớn. Một hệ thống có thể tổng hợp tốt khi được cấp đúng dữ liệu mà vẫn gặp khó khi phải tự tìm dữ liệu đó.

Vì vậy, cần phân biệt kết quả của phép thử này với trải nghiệm sử dụng trọn vẹn: từ câu hỏi ban đầu, tìm nguồn, xử lý, kiểm tra cho đến trả lời cuối cùng.

\subsection{Bộ chấm chưa kiểm mọi nội dung}
Bộ chấm chủ yếu xét những dữ kiện bắt buộc. Nếu câu trả lời có thêm lời khuyên hoặc nhận định ngoài danh sách đó, việc đạt đủ điểm không tự xác nhận các phần bổ sung đều đúng. Ngược lại, một ý đúng trong dữ liệu nhưng khác đáp án rút gọn cũng có thể bị chấm thiếu. \src{5}

Nghiên cứu tiếp theo cần có người đối chiếu một phần câu trả lời với nguồn, kiểm tra sự thống nhất giữa người và bộ chấm, đồng thời xem xét cả thông tin bị bỏ sót lẫn thông tin được thêm không có căn cứ.

\subsection{Chi phí và thời gian mới được đo trong phạm vi thử nghiệm}
Thời gian và token trong bảng kết quả thuộc quá trình sinh câu trả lời đã được ghi nhận. Chúng chưa phải toàn bộ chi phí triển khai, vận hành và chấm điểm. Thứ tự chạy hai cách cũng chưa được cân bằng ngẫu nhiên, nên khác biệt tải xử lý hoặc mạng có thể ảnh hưởng kết quả thời gian.

Tỷ lệ token lớn hơn cho thấy nhu cầu xử lý văn bản tăng, nhưng chưa đủ để chuyển trực tiếp thành một tỷ lệ tiền cụ thể nếu chưa xét loại mô hình, đầu vào, đầu ra và các thành phần liên quan.

\section{Hướng phát triển và đánh giá tiếp theo}

\subsection{Mở rộng và nâng chất lượng dữ liệu}
Trước hết cần tăng số nguồn kiến thức đáng tin cậy, cải thiện độ phủ tác phẩm và xác định rõ các ấn bản. Với thông tin biến động như giá hoặc tồn kho, nếu muốn phục vụ thực tế thì phải có nguồn cập nhật phù hợp và hiển thị rõ thời điểm thông tin được ghi nhận.

Nên tiếp tục kiểm tra các đánh giá thuộc đúng sản phẩm, phân biệt trải nghiệm giao hàng với chất lượng nội dung và giữ những trường hợp mâu thuẫn để hệ thống học cách diễn giải thận trọng. Mở rộng dữ liệu có căn cứ thường hữu ích hơn chỉ tăng số lượng tác tử trên cùng một tập thông tin nhỏ.

\subsection{Xây dựng bộ câu hỏi mới có độ khó thực chất}
Bộ câu hỏi tiếp theo nên có nhiều tình huống gần nhu cầu người dùng: nhiều tựa sách gần giống nhau, nhiều điều kiện kết hợp, thông tin thiếu, nguồn bất đồng, câu hỏi vượt phạm vi dữ liệu và yêu cầu giải thích vì sao chọn một cuốn thay vì cuốn khác.

Các câu hỏi dùng để kết luận phải được tách khỏi những câu đã dùng để chỉnh hệ thống. Nếu sửa cách xử lý sau khi xem một câu sai rồi tiếp tục dùng đúng câu đó để chứng minh hiệu quả, kết quả sẽ khó phản ánh khả năng xử lý tình huống mới.

\subsection{So sánh công bằng hai cách xử lý}
Cả một tác tử và nhiều tác tử cần được tiếp cận cùng loại thông tin, cùng yêu cầu bảo toàn tên sách, cùng tiêu chí trả lời và cùng quy tắc an toàn. Khi đó mới dễ xác định phần lợi ích do cách phối hợp, thay vì do một bên được hướng dẫn kỹ hơn.

Nên thay đổi cân bằng thứ tự chạy, lặp ở những thời điểm khác nhau và ghi đầy đủ kết quả. Các lần lặp của cùng một câu phải được xem là có liên quan, không coi tất cả là những bài toán mới để làm kết luận có vẻ chắc chắn hơn thực tế.

Một phần điều kiện cho đợt đánh giá tiếp theo đã được chuẩn bị. Trước khi chạy với mô hình thật, hệ thống ghi lại và khóa đúng phiên bản mã nguồn, từ chối chạy nếu mã nguồn có thay đổi chưa lưu, và yêu cầu người chạy đặt trần ngân sách. Bộ phiếu chấm mù cho hai người chấm độc lập cũng đã sẵn sàng, với mọi điểm để trống cho tới khi người thật điền. \src{9} Một lần chạy xác nhận với mô hình thật trên 60 hội thoại mới đã được thực hiện theo các điều kiện này và được trình bày ở phần kết quả; phần chấm của người thật vẫn chưa được thực hiện. \src{10}

\subsection{Đánh giá nhiều mặt thay vì chỉ một điểm số}
Một phép đánh giá đầy đủ nên trả lời các câu hỏi:
\begin{enumerate}
\item Câu trả lời đã đáp ứng đủ những ý người dùng hỏi chưa?
\item Các thông tin bắt buộc có đúng và đủ không?
\item Những nhận định bổ sung có được nguồn hỗ trợ không?
\item Hệ thống có biết hỏi lại hoặc từ chối đúng lúc không?
\item Người đọc có hiểu được lý do lựa chọn và giới hạn thông tin không?
\item Mức cải thiện có xứng đáng với thời gian chờ và tài nguyên tăng thêm không?
\end{enumerate}
Những câu hỏi này phù hợp hơn việc chỉ tuyên bố ``đạt 100\%'' mà không nói rõ điểm đó đang đo điều gì.

\section{Kết luận}
Hệ thống trợ lý tư vấn sách được xây dựng theo hướng phân chia công việc cho các tác tử chuyên trách, kết hợp thông tin sách, đánh giá người mua và nguồn kiến thức. Bộ phận điều phối tổ chức các phần việc, còn bước tổng hợp kiểm tra căn cứ trước khi đưa ra câu trả lời. Giá trị chính của cách tiếp cận là làm rõ trách nhiệm và nguồn thông tin của từng nhận định.

Trên bộ thực nghiệm 24 câu hỏi, nhiều tác tử đạt điểm trung bình 100\%, còn một tác tử đạt 98,96\%. Tuy nhiên, chênh lệch chỉ nằm ở một câu hỏi về hậu tố ấn bản, trong khi nhiều tác tử cần thời gian trung bình gấp 4,73 lần và lượng token gấp 2,78 lần. Vì vậy, kết quả hiện có cho thấy một đánh đổi cụ thể giữa mức đáp ứng dữ kiện và tài nguyên, chưa đủ chứng minh nhiều tác tử vượt trội trong mọi tình huống.

Hai đợt đánh giá lớn hơn với mô hình thật, mỗi đợt 720 lượt trên 60 hội thoại, cho cùng một bức tranh: không có bằng chứng nhiều tác tử hoàn thành tác vụ tốt hơn một tác tử. Ở lần chạy xác nhận có giả thuyết khai báo trước, nhiều tác tử không thua ở nhóm tác phẩm nào nhưng chênh lệch $+1{,}07$ điểm phần trăm chưa có ý nghĩa thống kê. Thành phần tạo khác biệt rõ là truy xuất có căn cứ: tắt truy xuất làm tỷ lệ hoàn thành giảm từ 16 đến 36 điểm phần trăm tùy đợt. Lần chạy xác nhận cũng cho thấy cả hai cách còn yếu khi câu hỏi được diễn đạt khác khuôn quen thuộc.

Sau thực nghiệm, nhóm đã sửa sáu điểm làm mất thông tin trong quá trình xử lý và loại khỏi kết luận các số liệu không truy ngược được tới kết quả chạy thật. Các sửa đổi được xác nhận bằng kiểm thử tự động; tác động của chúng lên chất lượng trả lời vẫn cần được đo lại.

Hướng phát triển phù hợp là tiếp tục nâng chất lượng dữ liệu, kiểm soát việc sử dụng bằng chứng và đánh giá trên các câu hỏi mới với điều kiện so sánh công bằng. Việc chọn một hay nhiều tác tử nên dựa trên nhu cầu thực tế và lợi ích đo được, thay vì mặc định số tác tử nhiều hơn sẽ cho kết quả tốt hơn.

\section*{Tài liệu tham khảo}
\begingroup
\small
\setlength{\parskip}{5pt}
\hypertarget{src-1}{}[1] \href{https://github.com/Thanh281105/multi-agent-system/blob/0a4930e87664dd91add96025a2230b92a15810e3/docs/architecture.md}{Tài liệu mô tả kiến trúc và phạm vi hoạt động của hệ thống.}\par

\hypertarget{src-2}{}[2] \href{https://github.com/Thanh281105/multi-agent-system/blob/0a4930e87664dd91add96025a2230b92a15810e3/docs/data.md}{Tài liệu nguồn dữ liệu sách và quy trình chuẩn bị dữ liệu.}\par

\hypertarget{src-3}{}[3] \href{https://github.com/Thanh281105/multi-agent-system/blob/0a4930e87664dd91add96025a2230b92a15810e3/data/knowledge/books-v1/coverage.md}{Phạm vi đối chiếu tác phẩm và giới hạn của nguồn kiến thức.}\par

\hypertarget{src-4}{}[4] \href{https://github.com/Thanh281105/multi-agent-system/blob/0a4930e87664dd91add96025a2230b92a15810e3/evaluation/v4_fair_benchmark_cases.json}{Bộ câu hỏi và các dữ kiện bắt buộc của thực nghiệm.}\par

\hypertarget{src-5}{}[5] \href{https://github.com/Thanh281105/multi-agent-system/blob/0a4930e87664dd91add96025a2230b92a15810e3/scripts/run_v4_full_benchmark.py}{Quy trình thực hiện và chấm điểm thực nghiệm.}\par

\hypertarget{src-6}{}[6] \href{https://github.com/Thanh281105/multi-agent-system/blob/0a4930e87664dd91add96025a2230b92a15810e3/output/evaluation-v4/full_24_benchmark_results.json}{Dữ liệu kết quả thực nghiệm được dùng để tính các bảng trong báo cáo.}\par

\hypertarget{src-7}{}[7] \href{https://github.com/Thanh281105/multi-agent-system/blob/0d83c799dbd8df2fdc5d43fb5a9e87125c3db33b/docs/benchmark-quality-p2.md}{Mô tả các điểm làm mất thông tin đã tái hiện, cách sửa và tám tình huống kiểm thử.}\par

\hypertarget{src-8}{}[8] \href{https://github.com/Thanh281105/multi-agent-system/blob/0d83c799dbd8df2fdc5d43fb5a9e87125c3db33b/docs/evidence-reporting.md}{Quy trình tạo báo cáo chỉ từ kết quả gốc và danh sách số liệu đã đính chính.}\par

\hypertarget{src-9}{}[9] \href{https://github.com/Thanh281105/multi-agent-system/blob/0d83c799dbd8df2fdc5d43fb5a9e87125c3db33b/docs/benchmark-readiness-p3.md}{Các điều kiện đã chuẩn bị cho đợt đánh giá tiếp theo với mô hình thật.}\par

\endgroup

\end{document}
